% vim: tw=78 ai spell spelllang=en_gb sw=2
% syntax spell toplevel

\documentclass[10pt]{llncs}
\usepackage{color}
\usepackage{url}
\usepackage[T1]{fontenc}
\usepackage{graphicx}
\usepackage{listings}
\usepackage{hyperref}
\usepackage{xspace}
\pagestyle{plain}

\lstset{language=C,tabsize=2,numbers=left,showstringspaces=false,
	breaklines=true,breakatwhitespace=true,escapeinside={(*@}{@*)},
	frame=l,captionpos=b,basicstyle=\sffamily\small}

\newcommand{\code}[2][]{\lstinline[columns=fullflexible, #1]$#2$}

\newcommand{\syz}{\textsc{SyzKaller}\xspace}
\newcommand{\trin}{\textsc{Trinity}\xspace}

\setcounter{tocdepth}{1}

\begin{document}

\title{SLE12-SP2 Kernel Testing}
\author{Jiri~Slaby}
\institute{Suse Labs \\
\email{jslaby@suse.cz}}

\maketitle

\begin{abstract}
  Throughout the past few years, many tools for testing the Linux kernel
  emerged. Some of them are really successful in finding bugs in the kernel.
  Since we use the kernel as a base in our products, we should also use the
  tools to find potential bugs in our products. There is an upcoming release
  of \textit{SUSE Linux Enterprise 12 Service Pack 2} (SLE12-SP2). And
  provided SLE12-SP2 happens to be rebased on a relatively new kernel version,
  it contains many yet uncovered bugs.

  This paper deals with the tools which were used by the author of this paper
  during the SLE12-SP2 kernel development. Apart from others, it especially
  covers \syz and \trin. It describes requirements, installation use and some
  of the found bugs in our kernels.
\end{abstract}

\section{Introduction} \label{sec:intro}
At SUSE, we develop our products based on Linux upstream code. The kernel is
no exception. Traditionally, we build our major releases of \textit{SUSE Linux
Enterprise} (SLE) on the current kernel versions. SLE10 is based on 2.6.16,
SLE11 on 2.6.27, SLE12 on 3.12, and so on. Sometimes we also upgrade to the
current kernel version on \textit{Service Pack} (SP) releases. For example
SLE11-SP1 contained 2.6.32. This happened also with SLE12-SP2. SLE12-SP2 will
be released with kernel 4.4.

Over the years, it was proven, that every such rebase onto a new version
brings regressions. Some drivers get removed over time because upstream
developers lost interest in fixing them. Contrary to our customers, of course.
Some functionality vanishes as it became obsolete. Not always for our
customers. Performance drops significantly sometimes as upstream does not focus
on performance primarily. And with new code, new bugs get introduced. While
the former two cases are intentional and can be easily reverted in most cases,
the latter two are not so easy to cover.

According to the \texttt{git} history between SLE12 (3.12) and SLE12-SP2 (4.4)
kernels, almost 40 thousand files were changed. More than seven million lines
of code were modified in the files\footnote{Measured by \texttt{git diff -w -b
-M -{}-shortstat -l 13000 v3.12..v4.4}.}. Therefore, with the manpower we have,
manual inspection of the changes is out of question. We have to rely on some
helpers to do the hard work. Performance issues are handled in most cases by
our \textit{Performance Team} and are not covered in this paper. We focus here
on tools aimed at finding functional bugs.

The range of the tools which can be used for kernel testing is quite large.
There are commercial ones like Coverity \textsc{Prevent} which is run on the
latest kernel by Coverity and results provided to
developers~\cite{coverity_scan}.  There are purely academic ones which are
run mostly by their developers to find bugs like \textsc{Stanse}~\cite{stanse}
or \textsc{LDV}~\cite{ldv}. Of the tools often used by kernel developers, the most popular occurs to be \textsc{Sparse}~\cite{sparse}.

In this paper, we describe tools used by the author of the paper for checking
the SLE12-SP2 kernel. He concentrated especially on tools which are not used
in day-to-day life of developers, but still they are relatively easy-to-use
and return good results. This includes \trin~\cite{trin} described in
Section~\ref{sec:trin}, \syz~\cite{syz} described in Section~\ref{sec:syz}. The
overall view, explanation of approaches and techniques, and comprehensive list
of tools can be found in the author's thesis~\cite{slaby2013automatic}.

\section{Fuzz Testing} \label{sec:fuzz}
Nowadays, one of the common techniques for exercising code is \emph{Fuzz
Testing}. In short, random data are generated and sent to program input.  The
reason behind is to confuse and potentially crash the program. The idea is
based on the assumption that programs do not handle their input correctly and
unhandled corner cases can cause havoc when triggered in the program. The
assumption is correct and fuzzers in this section support that.

\subsection{\trin} \label{sec:trin}
\cite{trin}

\paragraph{Usage} \trin is not a part of SLE. But it can be easily obtained
from the buildservice from project \texttt{devel:tools}. Since SLE12-SP2 is
not released yet, there is no SLE12-SP2 repository in there. But SLE-12-SP1
works nicely. Finally, \trin is not recommended to be run as superuser. So the
steps to install and run \trin are as follows:
\begin{lstlisting}[language=bash]
zypper ar -f http://download.opensuse.org/repositories/ \
             devel:/tools/SLE_12_SP1/devel:tools.repo
zypper install trinity
su some_user
trinity
\end{lstlisting}

\trin will run until a warning appears in the kernel log.

\subsection{\syz} \label{sec:syz}
\cite{syz}

\section{insmod \& rmmod} \label{sec:insmod}

\section{Results} \label{sec:res}

\section{Conclusion} \label{sec:concl}

%\section*{Availability}

\bibliographystyle{plain}
\bibliography{paper}

\end{document}
